% ModuloCounter_All.tex - Combined source listing for Lab_1_Modulo_n_Counter
\documentclass[11pt]{article}
\usepackage[utf8]{inputenc}
\usepackage[a4paper,margin=1in]{geometry}
\usepackage{xcolor}
\usepackage{courier}
\usepackage{listings}
\usepackage{hyperref}

% Color scheme (matches other labs)
\definecolor{kwcolor}{RGB}{0,0,180}
\definecolor{commentcolor}{RGB}{0,120,0}
\definecolor{stringcolor}{RGB}{180,0,0}
\definecolor{bg}{RGB}{248,248,248}

\lstset{
  backgroundcolor=\color{bg},
  language=C++,
  basicstyle=\footnotesize\ttfamily,
  keywordstyle=\color{kwcolor}\bfseries,
  commentstyle=\color{commentcolor}\itshape,
  stringstyle=\color{stringcolor},
  numberstyle=\tiny\color{gray},
  numbers=left,
  stepnumber=1,
  numbersep=8pt,
  tabsize=4,
  breaklines=true,
  showstringspaces=false,
  frame=single,
  captionpos=b
}

\begin{document}

\begin{center}
  {\LARGE Modulo N Counter - Combined Source Listing}
\end{center}
\vspace{6pt}

\section*{ModuloNDigit.h}
\begin{lstlisting}
/* ModuloNDigit.h - short notes and quick exam hints */

// ModuloNDigit: single digit counter that wraps at `rangeN`.
// Syntax note: member initialization and default values (C++11): `int countValue = 0;`.
// Exam hint: incrementCounter demonstrates modular arithmetic and wrap-around logic.
class ModuloNDigit
{
public:
    int countValue = 0;
    int rangeN;

    // Constructors/Destructors: initialize and cleanup
    ModuloNDigit();
    ModuloNDigit(int maxRange); // set upper bound for wrap-around
    ~ModuloNDigit();

    // incrementCounter: increases value and returns true if wrapped to 0
    bool incrementCounter();
    // printCounter: prints current value (logic kept simple for exams)
    void printCounter();
};
\end{lstlisting}

\section*{ModuloNDigit.cpp}
\begin{lstlisting}
/* ModuloNDigit.cpp - implementation with short notes */

#include "ModuloNDigit.h"
#include <iostream>
using namespace std;

// Constructor to initialize counter with a maximum range
// Constructor/Destructor: show basic initialization and cleanup
ModuloNDigit::ModuloNDigit(int maxRange)
{
    rangeN = maxRange; // store the modulo base
    countValue = 0; // ensure starting at 0
}

// Destructor
ModuloNDigit::~ModuloNDigit()
{
}

// Increment counter and reset if it reaches the maximum
// Logic: increment then check for wrap-around (modular arithmetic)
bool ModuloNDigit::incrementCounter()
{
    countValue++; // increase this digit
    if (countValue >= rangeN) // if reached base, wrap to 0
    {
        countValue = 0; // wrap-around
        return true; // indicate wrap occurred (used by ModuloNCounter to propagate carry)
    }
    return false; // no wrap
}

// Print the current value of the counter
// Function purpose: simple output for debugging or demonstration
void ModuloNDigit::printCounter()
{
    cout << countValue << endl;
}

// Default constructor: keep range at 0 until set by parameterized ctor
ModuloNDigit::ModuloNDigit()
{
    countValue = 0;
    rangeN = 0;
}
\end{lstlisting}

\section*{ModuloNCounter.h}
\begin{lstlisting}
/* ModuloNCounter.h - simple notes and exam hints */

// ModuloNCounter: manages a multi-digit counter (OOP: composition)
// - Uses an array of `ModuloNDigit` objects to represent digits
// Exam notes: constructors allocate resources; copy ctor must deep-copy; destructor frees memory.
class ModuloNCounter
{
public:
    ModuloNCounter(int maxRange, int noOfDigits);
    ModuloNCounter(const ModuloNCounter &other);
    ~ModuloNCounter(); // Destructor: frees dynamic array (see cpp)

    int rangeN; // base/range for each digit (syntax: simple member)
    int noOfDigits; // number of digits in the counter
    ModuloNDigit *counter; // pointer to dynamic array (pointers: holds address returned by new)

    // increment: advances the counter; returns true on overall overflow
    bool increment();
    // print: outputs current counter representation
    void print();
};
\end{lstlisting}

\section*{ModuloNCounter.cpp}
\begin{lstlisting}
/* ModuloNCounter.cpp - implementation with exam-friendly notes */

#include "ModuloNCounter.h"
#include "ModuloNDigit.h"
#include <iostream>

using namespace std;

// Character map for printing digits in bases up to 16 (strings represent digits)
// Syntax note: this is a C-style string/array used for mapping numeric value -> character.
char ch[] = "0123456789ABCDEF";

// Constructor: initializes counter array with specified range and digits
// OOP note: allocates dynamic array of `ModuloNDigit` (composition via pointer)
ModuloNCounter::ModuloNCounter(int maxRange, int noOfDigits)
{
    rangeN = maxRange; // store base/range
    this->noOfDigits = noOfDigits; // store digit count
    counter = new ModuloNDigit[noOfDigits]; // allocate dynamic array (remember to delete[] in destructor)
    for (int i = 0; i < noOfDigits; i++)
    {
        counter[i] = ModuloNDigit(rangeN); // initialize each digit with same range
    }
}

// Destructor: releases allocated memory for `counter`
// Constructors/Destructors: manage resource lifecycle (RAII note: this class uses manual new/delete)
ModuloNCounter::~ModuloNCounter()
{
    delete[] counter; // free dynamic memory to avoid leaks
}

// Increment the counter, checking each digit for wrap-around.
// Loops & conditions: iterate from least-significant to most-significant and handle carry-like behavior.
bool ModuloNCounter::increment()
{
    // Start from rightmost digit (index noOfDigits-1) and move left
    for (int index = noOfDigits - 1; index >= 0; index--)
    {
        // incrementCounter() returns true when this digit wrapped to 0 (like a carry)
        if (counter[index].incrementCounter())
        {
            if (index == 0) // leftmost digit wrapped => overall overflow
            {
                return true; // signal overflow to caller
            }
            // else continue loop to propagate carry to next digit on the left
        }
        else
        {
            break; // no wrap-around, stop propagation
        }
    }
    return false; // no overall overflow
}

// Copy constructor: performs deep copy to own a separate dynamic array
// Pointers/References note: copying pointer values would cause double-free; allocate new array instead.
ModuloNCounter::ModuloNCounter(const ModuloNCounter &other)
{
    this->rangeN = other.rangeN; // copy primitive member
    this->noOfDigits = other.noOfDigits; // copy count
    counter = new ModuloNDigit[noOfDigits]; // allocate fresh array
    for (int i = 0; i < noOfDigits; i++)
    {
        counter[i] = other.counter[i]; // copy each ModuloNDigit (calls its assignment operator)
    }
}

// Print current counter value using `ch` mapping (no STL used here)
// Function purpose: produce human-readable representation in configured base.
void ModuloNCounter::print()
{
    for (int i = 0; i < noOfDigits; i++)
    {
        // Access digit value and map to character (array indexing and member access)
        cout << ch[counter[i].countValue];
    }
    cout << " "; // spacing between printed values
}
\end{lstlisting}

\section*{main.cpp}
\begin{lstlisting}
/* main.cpp - simple runner with exam-friendly inline notes */

#include <iostream>
#include <cstdlib>
#include "ModuloNCounter.h"

using namespace std;

int main()
{
    // `setvbuf` used to disable stdout buffering for immediate console output
    setvbuf(stdout, NULL, _IONBF, 0);
    int noOfDigits, typeOfCounter;

    // Main interactive loop - demonstrates program flow, input and loops
    while (true)
    {
        cout << endl << endl << "Please enter the parameters of your counter:" << endl;

        cout << "number of digits: ";
        cin >> noOfDigits; // user input (no validation here)

        cout << "type (2/8/10/16): ";
        cin >> typeOfCounter; // base of counter (syntax: integer variable)

        // Print human-readable description of chosen base
        cout << noOfDigits << " digit ";
        if (typeOfCounter == 2)
            cout << "binary counter";
        else if (typeOfCounter == 8)
            cout << "octal counter";
        else if (typeOfCounter == 10)
            cout << "decimal counter";
        else if (typeOfCounter == 16)
            cout << "hexadecimal counter";
        cout << endl;

        // Initialize the counter object (calls constructor)
        ModuloNCounter myCounter = ModuloNCounter(typeOfCounter, noOfDigits);

        // Compute number of states: base^digits using a loop (avoid pow to keep it simple)
        int maxCounterValue = 1;
        for (int i = 0; i < noOfDigits; i++)
        {
            // Loop and arithmetic: repeated multiplication produces base^digits
            maxCounterValue *= typeOfCounter;
        }

        // Iterate and print counter values; note: +2 shows wrap-around examples
        for (int i = 0; i < maxCounterValue + 2; i++)
        {
            if (i % typeOfCounter == 0)
                cout << endl; // formatting: new line every `typeOfCounter` prints
            myCounter.print(); // prints current value
            myCounter.increment(); // advance counter (may cause wrap)
        }

    }

    return 0; // program never reaches here due to infinite loop
}
\end{lstlisting}

\end{document}
